%
% File main.tex
%
% NLP course editorial-tips document, Warsaw University of Technology.
% Prepared by Anna Wróblewska, adapted from the ACL 2015 conference
% formatting template (acl2015.sty), originally by the ACL 2015
% organizers and earlier ACL/EACL format authors (see acl2015.pdf
% for the full credits chain).

\documentclass[11pt]{article}
\usepackage{acl2015}
\usepackage{times}
\usepackage{url}
\usepackage{latexsym}
\usepackage{xcolor}

%\setlength\titlebox{5cm}

% You can expand the title box if you need extra space
% to show all the authors. Please do not make the title box
% smaller than 5cm (the original size); we will check this
% in the camera-ready version and ask you to change it back.


\title{Editorial Tips\\NLP Course, Winter 2026/27}

\author{Author 1 (submitted by) \\
  Album number \\
  {\tt email@pw.edu.pl} \\\And
  Author 2 \\
  Album number \\
  {\tt email@pw.edu.pl} \\ \And
  supervisor: Anna Wróblewska\\
  Warsaw University of Technology \\
  {\tt anna.wroblewska1@pw.edu.pl}}

\date{}

\begin{document}
\maketitle

{\small\em Template: Anna Wr\'oblewska (Warsaw University of Technology), adapted from the ACL 2015 conference formatting template.}
\vspace{4pt}

\begin{abstract}
  This document collects the formatting and editorial rules for NLP-course reports. Read it once, then use it as a checklist before every submission.
\end{abstract}

\section{What stays the same across all phases}

\subsection{Layout and fonts}
Two columns, 11pt Times (already set by \texttt{acl2015.sty} and \texttt{\textbackslash usepackage\{times\}} above -- do not change the class options). Keep section headings numbered (Arabic numerals); do not number subsubsections. A4 page size is handled by the style file.

\subsection{Citations}
Use \verb|\cite{key}| for parenthetical citations and \verb|\shortcite{key}| or \verb|\newcite{key}| for ``Author (Year)'' in running text, exactly as in the original ACL macros. Two graded rules from the criteria document govern \emph{which} papers you cite, not just how:
\begin{itemize}
\item Preferred venues: ACL, EMNLP, NAACL, or ACM Computing Surveys (aclanthology.org).
\item Cite the published version of a work whenever one exists; preprint-only citations (arXiv with no venue) should be under one third of your bibliography. The final-report template additionally excludes arXiv preprints from ``well-justified literature'' unless the work is too recent to have a published version yet.
\end{itemize}

\subsection{Language}
Scientific, precise language; no unexplained abbreviations on first use. Every figure and table caption must be understandable without reading the surrounding text -- a reader skimming only captions should still follow the story.

\subsection{Reproducibility}
Seeds, library versions, and the exact command(s) used to produce your main reported result must be stated somewhere in the report or its appendix. Use \texttt{reproducibility-checklist.tex} as the appendix template -- it is graded as part of Phase 3.

\section{What differs by phase}

\begin{itemize}
\item \textbf{Phase 1 (proposal)}: max. 4 meaningful pages. Template: \texttt{project-proposal.tex}.
\item \textbf{Phase 3 (final report)}: about 9 pages, max 12, plus the reproducibility appendix. Template: \texttt{project-final-report-structure.tex}.
\end{itemize}
Do not carry one phase's page limit over to another -- this has been a recurring point of confusion.

\section{Mandatory content, every phase}

Three items are conditions of acceptance, not scored criteria: a report missing any of them is not graded at all until it is added, in every phase where it applies.
\begin{itemize}
\item \textbf{Contribution table} -- each team member, what they did, workload in hours. Required from Phase 2 onward.
\item \textbf{Declaration of LLM use} -- which tools, for what.
\item \textbf{Data statement} -- source of the data, licence, whether it contains personal data, and what you do about it.
\end{itemize}

\section{Submission mechanics}

\begin{itemize}
\item \textbf{Folder naming}: every milestone goes in a folder named \texttt{surname\_surname\_surname\_surname} (all team members). The order set at the first submission must not change for the rest of the semester -- a mismatch costs 5 points.
\item \textbf{Format}: PDF compiled from these LaTeX templates, or the Word equivalent if you were told to use one. \textbf{Wrong template: $-5$ points}, and you resubmit in the correct one.
\item \textbf{Late submission}: one delay of up to 48 hours per semester carries no penalty. Every further delay costs 5 points per started week of that phase.
\end{itemize}

\end{document}
